% !TEX spellcheck = en_US



% ~~~~~~~~~~~~~~~~~~~~~~~~~~~~~~~~~~~~~~~ V
% (NC) 2026 Haluk Bingol
% github.com/halukbingol/LaTeX-Templates
%
% Licensees may copy, distribute, display, and perform the work and make derivative works 
% and remixes based on it only for non-commercial purposes. 
% ~~~~~~~~~~~~~~~~~~~~~~~~~~~~~~~~~~~~~~~ A




% =====================================================================
%  Java on the Command Line -- Freshman Level: Your First Programs in a Terminal
%  ---------------------------------------------------------------------
%  Built on hbCisLaTeXTemplate.tex with the libraries in ../../hblib/.
%  Programs and their outputs are read from codeoutput/:
%     codeoutput/<Name>.java   the program
%     codeoutput/<Name>.in     keyboard input (optional)
%     codeoutput/<Name>.txt    what it printed, made by:  sh run_examples.sh
% =====================================================================

\documentclass[aspectratio=169,10pt,compress]{beamer}

% ~~~~~~~~~~~~~~~~~~~~~~~~~~~~~~~~~~~~~~~~~~~~~~~~~~~~~~~~~~~~~~~~~ V hblib
	\usepackage{../../hblib/hblibPresentation} % lib presentation: theme, i/N, \hSection, algorithm2e, ...
	\usepackage{../../hblib/hblibCodeoutput} % lib code + output: \lstcode, \lstcodedown, ...
	\usepackage{../../hblib/hblibDatastructures} % lib data structures: \hString, \hStack, \hQueue
% ~~~~~~~~~~~~~~~~~~~~~~~~~~~~~~~~~~~~~~~~~~~~~~~~~~~~~~~~~~~~~~~~~ A hblib


% xcolor, array (and the L{width} column), listings, tikz and algorithm2e come from the hblib libraries
\usepackage[T1]{fontenc}
\usepackage{lmodern}
\usepackage{booktabs}

\definecolor{gitred}{RGB}{240,80,50}
\newcommand{\cmd}[1]{\texttt{\color{gitred}#1}}

\tikzset{
  box/.style={draw, thick, rounded corners, align=center, font=\small, inner sep=5pt},
  tool/.style={-{Stealth[length=2.5mm]}, thick}
}

% ---------------------------------------------------------------------
\title{Java on the Command Line}
\subtitle{Freshman Level: Your First Programs in a Terminal}
\author[Bingol]{Haluk O. Bingol}

\institute[FCIS]{
    Faculty of Computer and Information Sciences\\
    Yeditepe University
}
\date{
	{\tiny \hbTimeStamp}
}

\begin{document}




% ~~~~~~~~~~~~~~~~~~~~~~~~~~~~~~~~~~~~~~ V frm
\begin{frame}[t,fragile]
  \titlepage
\end{frame}
% ~~~~~~~~~~~~~~~~~~~~~~~~~~~~~~~~~~~~~~ A frm




% ~~~~~~~~~~~~~~~~~~~~~~~~~~~~~~~~~~~~~~ V frm
\begin{frame}[t,fragile] \frametitle{Outline}
  \begin{columns}[T]
    \begin{column}{0.48\textwidth}
      \tableofcontents[sections={1-3}, hideallsubsections]
    \end{column}
    \begin{column}{0.48\textwidth}
      \tableofcontents[sections={4-6}, hideallsubsections]
    \end{column}
  \end{columns}
\end{frame}
% ~~~~~~~~~~~~~~~~~~~~~~~~~~~~~~~~~~~~~~ A frm

% ~~~~~~~~~~~~~~~~~~~~~~~~~~~~~~~~~~~~~~ V
\hSection[Terminal]{The Terminal}{}{}
% ~~~~~~~~~~~~~~~~~~~~~~~~~~~~~~~~~~~~~~ A

% ~~~~~~~~~~~~~~~~~~~~~~~~~~~~~~~~~~~~~~ V
\subsection[Open]{Opening a terminal}
% ~~~~~~~~~~~~~~~~~~~~~~~~~~~~~~~~~~~~~~ A




% ~~~~~~~~~~~~~~~~~~~~~~~~~~~~~~~~~~~~~~ V frm
\begin{frame}[t,fragile] \frametitle{Opening a terminal}
  \begin{columns}[T]
    \begin{column}{0.48\textwidth}
      \begin{block}{Windows}
        \small \texttt{Win} key, type \hCode{cmd} (Command Prompt) or \hCode{terminal} (Windows Terminal), Enter.
        In File Explorer: type \hCode{cmd} in the address bar to open it \textbf{in that folder}.
      \end{block}
    \end{column}
    \begin{column}{0.48\textwidth}
      \begin{block}{macOS}
        \small \texttt{Cmd}+\texttt{Space}, type \hCode{Terminal}, Enter.
        In Finder: Control-click a folder, \emph{Services}, \emph{New Terminal at Folder} (turn it on in
        \emph{System Settings, Keyboard, Keyboard Shortcuts, Services} if missing).
      \end{block}
    \end{column}
  \end{columns}
  \vspace{0.3em}
  \footnotesize
  \begin{tabular}{@{}L{0.3\textwidth}L{0.3\textwidth}L{0.3\textwidth}@{}}
    \toprule
    \textbf{To \dots}                 & \textbf{Windows (cmd)}        & \textbf{macOS (zsh)} \\
    \midrule
    see where you are                 & \hCode{cd}                   & \hCode{pwd} \\
    list files                        & \hCode{dir}                  & \hCode{ls} \\
    go into a folder / up             & \hCode{cd week1} / \hCode{cd ..} & \hCode{cd week1} / \hCode{cd ..} \\
    make a folder                     & \hCode{mkdir week1}          & \hCode{mkdir week1} \\
    show a file                       & \hCode{type Hello.java}      & \hCode{cat Hello.java} \\
    clear the screen                  & \hCode{cls}                  & \hCode{clear} \\
    \bottomrule
  \end{tabular}
\end{frame}
% ~~~~~~~~~~~~~~~~~~~~~~~~~~~~~~~~~~~~~~ A frm \hCode




% ~~~~~~~~~~~~~~~~~~~~~~~~~~~~~~~~~~~~~~ V
\subsection[Keys]{Useful keys}
% ~~~~~~~~~~~~~~~~~~~~~~~~~~~~~~~~~~~~~~ A




% ~~~~~~~~~~~~~~~~~~~~~~~~~~~~~~~~~~~~~~ V frm
\begin{frame}[t,fragile] \frametitle{Keys that save time}
  \small
  \begin{tabular}{@{}L{0.25\textwidth}L{0.65\textwidth}@{}}
    \toprule
    \textbf{Key}                      & \textbf{Does} \\
    \midrule
    \texttt{Tab}                      & completes a file or folder name: \hCode{javac He}\texttt{Tab} $\to$ \hCode{javac Hello.java} \\
    \texttt{Up} / \texttt{Down}       & previous / next command (no retyping after an edit) \\
    \texttt{Ctrl}+\texttt{C}          & stops a running program (e.g.\ an endless loop) \\
    \texttt{Ctrl}+\texttt{D} (macOS), \texttt{Ctrl}+\texttt{Z} Enter (Windows) & ends keyboard input (``end of file'') \\
    \texttt{Ctrl}+\texttt{L} (macOS), \texttt{cls} (Windows) & clears the screen \\
    \bottomrule
  \end{tabular}

  \vspace{0.6em}
  \begin{alertblock}{Spaces in folder names}
    \small \hCode{cd My Documents} fails; write \hCode{cd "My Documents"}. Simplest: use folder names without spaces or Turkish letters,
    e.g.\ \hCode{cse101/week1}.
  \end{alertblock}
\end{frame}
% ~~~~~~~~~~~~~~~~~~~~~~~~~~~~~~~~~~~~~~ A frm

% ~~~~~~~~~~~~~~~~~~~~~~~~~~~~~~~~~~~~~~ V
\hSection[Install]{Installing and Checking Java}{}{}
% ~~~~~~~~~~~~~~~~~~~~~~~~~~~~~~~~~~~~~~ A

% ~~~~~~~~~~~~~~~~~~~~~~~~~~~~~~~~~~~~~~ V
\subsection[Install]{Installing}
% ~~~~~~~~~~~~~~~~~~~~~~~~~~~~~~~~~~~~~~ A




% ~~~~~~~~~~~~~~~~~~~~~~~~~~~~~~~~~~~~~~ V frm
\begin{frame}[t,fragile] \frametitle{Install a JDK}
  \small Get \textbf{Eclipse Temurin 21 (LTS)} from \texttt{adoptium.net}: a free JDK, which contains \texttt{javac}, \texttt{java} and \texttt{jar}.
  \begin{columns}[T]
    \begin{column}{0.48\textwidth}
      \begin{block}{Windows}
        \footnotesize
        \begin{enumerate}
          \item Download the \texttt{.msi} for \emph{Windows x64}
          \item In the installer, enable \emph{Add to PATH} and \emph{Set JAVA\_HOME}
          \item Or in a terminal: \cmd{winget install EclipseAdoptium.Temurin.21.JDK}
          \item Close and \textbf{reopen} the terminal
        \end{enumerate}
      \end{block}
    \end{column}
    \begin{column}{0.48\textwidth}
      \begin{block}{macOS}
        \footnotesize
        \begin{enumerate}
          \item Download the \texttt{.pkg}: \emph{aarch64} for Apple silicon (M1, M2, \dots), \emph{x64} for Intel
          \item Double-click and follow the installer
          \item Or with Homebrew: \hCode{brew install --cask temurin@21}
          \item Open a \textbf{new} Terminal window
        \end{enumerate}
      \end{block}
    \end{column}
  \end{columns}
  \small Check in the new terminal: \hCode{java -version} and \hCode{javac -version} must both print \texttt{21...}.
\end{frame}
% ~~~~~~~~~~~~~~~~~~~~~~~~~~~~~~~~~~~~~~ A frm




% ~~~~~~~~~~~~~~~~~~~~~~~~~~~~~~~~~~~~~~ V
\subsection[PATH]{When the command is not found}
% ~~~~~~~~~~~~~~~~~~~~~~~~~~~~~~~~~~~~~~ A




% ~~~~~~~~~~~~~~~~~~~~~~~~~~~~~~~~~~~~~~ V frm
\begin{frame}[t,fragile] \frametitle{``Command not found}
  \small If the shell cannot find \texttt{javac} on the \textbf{PATH}, you see one of these:


% +++++++++++++++++++++++++++++++++++++++ V lst
%: -Lst.
\begin{lstlisting}[style=codesnippet, language={}, basicstyle=\ttfamily\scriptsize]
C:\Users\ayse> javac -version
'javac' is not recognized as an internal or external command,
operable program or batch file.
\end{lstlisting}
% +++++++++++++++++++++++++++++++++++++++ A lst




% +++++++++++++++++++++++++++++++++++++++ V lst
%: -Lst.
\begin{lstlisting}[style=codesnippet, language={}, basicstyle=\ttfamily\scriptsize]
ayse@mac ~ % javac -version
The operation couldn't be completed. Unable to locate a Java Runtime that supports javac.
Please visit http://www.java.com for information on installing Java.
\end{lstlisting}
% +++++++++++++++++++++++++++++++++++++++ A lst



  \begin{columns}[T]
    \begin{column}{0.48\textwidth}
      \begin{block}{Windows fix}
        \footnotesize Reopen the terminal. Still failing: \emph{Edit the system environment variables} $\to$ \emph{Path} $\to$
        add \texttt{C:\textbackslash Program Files\textbackslash Eclipse Adoptium\textbackslash jdk-21...\textbackslash bin}. Check with \hCode{where javac}.
      \end{block}
    \end{column}
    \begin{column}{0.48\textwidth}
      \begin{block}{macOS fix}
        \footnotesize macOS has only placeholder \texttt{java}/\texttt{javac}: this message means no JDK is installed yet.
        List installed JDKs: \cmd{/usr/libexec/java\_home -V}; check \hCode{which javac}.
      \end{block}
    \end{column}
  \end{columns}
\end{frame}
% ~~~~~~~~~~~~~~~~~~~~~~~~~~~~~~~~~~~~~~ A frm




% ~~~~~~~~~~~~~~~~~~~~~~~~~~~~~~~~~~~~~~ V
\hSection[First]{Your First Program}{}{}
% ~~~~~~~~~~~~~~~~~~~~~~~~~~~~~~~~~~~~~~ A




% ~~~~~~~~~~~~~~~~~~~~~~~~~~~~~~~~~~~~~~ V
\subsection[Cycle]{Edit, compile, run}
% ~~~~~~~~~~~~~~~~~~~~~~~~~~~~~~~~~~~~~~ A




% ~~~~~~~~~~~~~~~~~~~~~~~~~~~~~~~~~~~~~~ V frm
\begin{frame}[t,fragile] \frametitle{Edit, compile, run}
  \begin{center}
  \begin{tikzpicture}[x=1cm, y=1cm]
    \node[box, fill=blue!8] (e) at (0,0) {\textbf{edit}\\ \scriptsize\texttt{Hello.java}\\ \scriptsize in an editor};
    \node[box, fill=orange!12] (c) at (4.5,0) {\textbf{compile}\\ \scriptsize\texttt{javac Hello.java}};
    \node[box, fill=green!12] (r) at (9,0) {\textbf{run}\\ \scriptsize\texttt{java Hello}};
    \draw[tool] (e) -- (c);
    \draw[tool] (c) -- node[above, font=\scriptsize]{no errors} (r);
    \draw[tool, red!70!black] (c) to[bend left=35] node[below, font=\scriptsize]{compile errors} (e);
    \draw[tool, red!70!black] (r) to[bend left=40] node[below, font=\scriptsize]{wrong output or exception} (e);
  \end{tikzpicture}
  \end{center}
  \begin{itemize}\small
    \item Use a plain-text editor: VS Code, Notepad++ or Notepad (Windows), TextEdit in \emph{plain text} mode (macOS); save as \textbf{UTF-8}
    \item Keep the terminal open in the same folder as your file
    \item After every change: \textbf{save}, compile again, run again
  \end{itemize}
\end{frame}
% ~~~~~~~~~~~~~~~~~~~~~~~~~~~~~~~~~~~~~~ A frm

% ~~~~~~~~~~~~~~~~~~~~~~~~~~~~~~~~~~~~~~ V
\subsection[First]{A first session}
% ~~~~~~~~~~~~~~~~~~~~~~~~~~~~~~~~~~~~~~ A




% ~~~~~~~~~~~~~~~~~~~~~~~~~~~~~~~~~~~~~~ V frm
\begin{frame}[t,fragile] \frametitle{A first session}
  \lstcode[basicstyle=\ttfamily\scriptsize]{codeoutput/proj/First/Hello.java}{Java}
  \uncover<2->{\lstoutput[basicstyle=\ttfamily\tiny]{codeoutput/First.txt}}

  \vspace{0.2em}
  \small Windows: \hCode{dir} instead of \hCode{ls}, \hCode{type} instead of \hCode{cat}; \hCode{javac} and \hCode{java} are the same.
\end{frame}
% ~~~~~~~~~~~~~~~~~~~~~~~~~~~~~~~~~~~~~~ A frm

% ~~~~~~~~~~~~~~~~~~~~~~~~~~~~~~~~~~~~~~ V
\hSection[Errors]{Reading Error Messages}{}{}
% ~~~~~~~~~~~~~~~~~~~~~~~~~~~~~~~~~~~~~~ A

% ~~~~~~~~~~~~~~~~~~~~~~~~~~~~~~~~~~~~~~ V
\subsection[Compile]{Compile-time errors}
% ~~~~~~~~~~~~~~~~~~~~~~~~~~~~~~~~~~~~~~ A




% ~~~~~~~~~~~~~~~~~~~~~~~~~~~~~~~~~~~~~~ V frm
\begin{frame}[t,fragile] \frametitle{The file name must match the class}
  \lstcode[basicstyle=\ttfamily\scriptsize]{codeoutput/proj/NameMismatch/Greeting.java}{Java}
  \uncover<2->{\lstoutput[basicstyle=\ttfamily\tiny]{codeoutput/NameMismatch.txt}}

  \vspace{0.2em}
  \small A \hCode{public class Hello} must be in \texttt{Hello.java} (capital letters matter). Windows: \hCode{ren} instead of \hCode{mv}.
\end{frame}
% ~~~~~~~~~~~~~~~~~~~~~~~~~~~~~~~~~~~~~~ A frm




% ~~~~~~~~~~~~~~~~~~~~~~~~~~~~~~~~~~~~~~ V frm
\begin{frame}[t,fragile] \frametitle{A missing semicolon}
  \lstcode[basicstyle=\ttfamily\scriptsize]{codeoutput/proj/Semicolon/Hello.java}{Java}
  \uncover<2->{\lstoutput[basicstyle=\ttfamily\tiny]{codeoutput/Semicolon.txt}}

  \vspace{0.2em}
  \small \hRemark{Read the message:} file and line (\texttt{Hello.java:3}), 
  what is wrong (\texttt{';' expected}),
  and a \texttt{\textasciicircum} under the place. 
  Fix the \textbf{first} error first; later ones may disappear.
\end{frame}
% ~~~~~~~~~~~~~~~~~~~~~~~~~~~~~~~~~~~~~~ A frm




% ~~~~~~~~~~~~~~~~~~~~~~~~~~~~~~~~~~~~~~ V frm
\begin{frame}[t,fragile] \frametitle{A wrong capital letter}
  \lstcode[basicstyle=\ttfamily\scriptsize]{codeoutput/proj/Typo/Hello.java}{Java}
  \uncover<2->{\lstoutput[basicstyle=\ttfamily\tiny]{codeoutput/Typo.txt}}

  \vspace{0.2em}
  \small Java is case-sensitive: \hCode{system} is not \hCode{System}. The message can be indirect; look at the
  line and the \texttt{\textasciicircum}.
\end{frame}
% ~~~~~~~~~~~~~~~~~~~~~~~~~~~~~~~~~~~~~~ A frm

% ~~~~~~~~~~~~~~~~~~~~~~~~~~~~~~~~~~~~~~ V
\subsection[Run]{Run-time errors}
% ~~~~~~~~~~~~~~~~~~~~~~~~~~~~~~~~~~~~~~ A




% ~~~~~~~~~~~~~~~~~~~~~~~~~~~~~~~~~~~~~~ V frm
\begin{frame}[t,fragile] \frametitle{``Could not find or load main class}
  \lstoutput[basicstyle=\ttfamily\scriptsize]{codeoutput/RunErrors.txt}

  \vspace{0.2em}
  \small Give the \textbf{class name} (no \hCode{.class}), with the right capitals, from the \textbf{folder} that contains
  \hCode{Hello.class}. On Windows and macOS \hCode{java hello} 
  gives \texttt{NoClassDefFoundError: Hello (wrong name: hello)} instead,
  because their file systems ignore case.
\end{frame}
% ~~~~~~~~~~~~~~~~~~~~~~~~~~~~~~~~~~~~~~ A frm




% ~~~~~~~~~~~~~~~~~~~~~~~~~~~~~~~~~~~~~~ V frm
\begin{frame}[t,fragile] \frametitle{No \texttt{main} method}
  \lstcode[basicstyle=\ttfamily\scriptsize]{codeoutput/proj/NoMain/Hello.java}{Java}
  \uncover<2->{\lstoutput[basicstyle=\ttfamily\tiny]{codeoutput/NoMain.txt}}

  \vspace{0.2em}
  \small \hCode{javac} accepts the class; only \hCode{java} needs \hCode{public static void main(String[] args)}.
\end{frame}
% ~~~~~~~~~~~~~~~~~~~~~~~~~~~~~~~~~~~~~~ A frm




% ~~~~~~~~~~~~~~~~~~~~~~~~~~~~~~~~~~~~~~ V frm
\begin{frame}[t,fragile] \frametitle{An exception: the program crashes}
  \lstcode[basicstyle=\ttfamily\tiny]{codeoutput/proj/Divide/Divide.java}{Java}
  \uncover<2->{\lstoutput[basicstyle=\ttfamily\tiny]{codeoutput/Divide.txt}}

  \vspace{0.2em}
  \small A \textbf{stack trace}: the exception, its message, and where it happened (\hCode{Divide.java:5}). The exit code is 1.
\end{frame}
% ~~~~~~~~~~~~~~~~~~~~~~~~~~~~~~~~~~~~~~ A frm




% ~~~~~~~~~~~~~~~~~~~~~~~~~~~~~~~~~~~~~~ V
\hSection[IO]{Input and Arguments}{}{}
% ~~~~~~~~~~~~~~~~~~~~~~~~~~~~~~~~~~~~~~ A




% ~~~~~~~~~~~~~~~~~~~~~~~~~~~~~~~~~~~~~~ V
\subsection[Input]{Keyboard input}
% ~~~~~~~~~~~~~~~~~~~~~~~~~~~~~~~~~~~~~~ A




% ~~~~~~~~~~~~~~~~~~~~~~~~~~~~~~~~~~~~~~ V frm
\begin{frame}[t,fragile] \frametitle{Reading input}
  \lstcode[basicstyle=\ttfamily\scriptsize]{codeoutput/proj/Greet/Greet.java}{Java}
  \uncover<2->{\lstoutput[basicstyle=\ttfamily\scriptsize]{codeoutput/Greet.txt}}

  \vspace{0.2em}
  \small Run \hCode{java Greet} and type a name. For testing, \hCode{<} reads the input from a file and
  \hCode{|} (a \textbf{pipe}) passes another command's output; both work on Windows and macOS.
\end{frame}
% ~~~~~~~~~~~~~~~~~~~~~~~~~~~~~~~~~~~~~~ A frm




% ~~~~~~~~~~~~~~~~~~~~~~~~~~~~~~~~~~~~~~ V
\subsection[Args]{Arguments}
% ~~~~~~~~~~~~~~~~~~~~~~~~~~~~~~~~~~~~~~ A




% ~~~~~~~~~~~~~~~~~~~~~~~~~~~~~~~~~~~~~~ V frm
\begin{frame}[t,fragile] \frametitle{Arguments: \hCode{java Sum 3 4 5}}
  \lstcode[basicstyle=\ttfamily\scriptsize]{codeoutput/proj/Sum/Sum.java}{Java}
  \uncover<2->{\lstoutput[basicstyle=\ttfamily\tiny, linerange={1-7,11-12}]{codeoutput/Sum.txt}}

  \vspace{0.2em}
  \small Arguments arrive in \hCode{args} as strings; \hCode{Integer.parseInt} makes numbers (3 internal stack lines skipped).
\end{frame}
% ~~~~~~~~~~~~~~~~~~~~~~~~~~~~~~~~~~~~~~ A frm




% ~~~~~~~~~~~~~~~~~~~~~~~~~~~~~~~~~~~~~~ V
\hSection[Summary]{Summary}{}{}
% ~~~~~~~~~~~~~~~~~~~~~~~~~~~~~~~~~~~~~~ A




% ~~~~~~~~~~~~~~~~~~~~~~~~~~~~~~~~~~~~~~ V
\subsection[Errors]{Error message guide}
% ~~~~~~~~~~~~~~~~~~~~~~~~~~~~~~~~~~~~~~ A




% ~~~~~~~~~~~~~~~~~~~~~~~~~~~~~~~~~~~~~~ V frm
\begin{frame}[t,fragile] \frametitle{Error message guide}
  \footnotesize
  \begin{tabular}{@{}L{0.44\textwidth}L{0.5\textwidth}@{}}
    \toprule
    \textbf{Message}                                                   & \textbf{Usual cause} \\
    \midrule
    \hCode{'javac' is not recognized} / \hCode{command not found}    & JDK not installed or not on PATH; reopen the terminal \\
    \hCode{class X is public, should be declared in a file named X.java} & file name and class name differ \\
    \hCode{';' expected}, \hCode{cannot find symbol}                 & typo, missing semicolon, wrong capital letter \\
    \hCode{Could not find or load main class}                         & wrong folder, wrong name, or \texttt{.class} typed \\
    \hCode{Main method not found}                                     & \hCode{main} missing or not \hCode{public static void main(String[] args)} \\
    \hCode{Exception in thread "main" ...}                            & a run-time error; read the type, the message and the line \\
    \bottomrule
  \end{tabular}
\end{frame}
% ~~~~~~~~~~~~~~~~~~~~~~~~~~~~~~~~~~~~~~ A frm




% ~~~~~~~~~~~~~~~~~~~~~~~~~~~~~~~~~~~~~~ V
\subsection[Exercises]{Exercises}
% ~~~~~~~~~~~~~~~~~~~~~~~~~~~~~~~~~~~~~~ A




% ~~~~~~~~~~~~~~~~~~~~~~~~~~~~~~~~~~~~~~ V frm
\begin{frame}[t,fragile] \frametitle{Exercises}
  \begin{enumerate}\small
    \item In a folder \hCode{cse101/week1}, write, compile and run \hCode{Hello.java} on your own computer.
    \item Make each error from these slides on purpose; read the message before fixing it.
    \item Write \hCode{Max.java} that prints the largest of its integer arguments: \hCode{java Max 3 9 4} $\to$ \texttt{9}.
    \item Write \hCode{Echo.java} that prints every line it reads, numbered. Run it with the keyboard
          (end with \texttt{Ctrl}+\texttt{D} on macOS, \texttt{Ctrl}+\texttt{Z} Enter on Windows) and with \hCode{< Echo.java}.
    \item Run \hCode{java Hello.java} (one step). When is this enough, and when do you need \hCode{javac}?
    \item Help a friend with the other operating system: write the same steps for Windows \emph{and} macOS.
  \end{enumerate}
\end{frame}
% ~~~~~~~~~~~~~~~~~~~~~~~~~~~~~~~~~~~~~~ A frm




% ~~~~~~~~~~~~~~~~~~~~~~~~~~~~~~~~~~~~~~ V frm
\begin{frame}[t,fragile]
	\begin{center}
		\vspace{4cm}
		\hRemark{\Huge Questions?}
	\end{center}
\end{frame}
% ~~~~~~~~~~~~~~~~~~~~~~~~~~~~~~~~~~~~~~ A frm



\end{document}
